\begin{helper}[First normal form in the last-index case]
Let \((\varphi,T)\) be a minimum bad pair with
\(2<\rho(T)=|E(T)|=n\).  There are a coloring \(\varphi_1\), a color
\(\beta\), and an index \(j\le n-2\) such that
\((\varphi_1,T)\) is a minimum bad pair, \(\beta\) is missing at both
\(y_j,y_n\), and \(\beta\) is not used on any colored displayed tree
edge.
\end{helper}
\begin{proof}
The inequality gives \(n\ge3\), so the prefixes through \(y_{n-1}\) and
\(y_{n-2}\) are defined by \noderef{TreeSequencePrefix}.  Unfolding
\noderef{TashTreeMinimumBadPair} and then \noderef{TashTreeBadPair}, the
original pair supplies a \(k\)-triple, the \noderef{TashkinovTree}
condition on \(T\), and non-elementariness of the full vertex set.  In
particular, \noderef{KTriple} supplies properness of \(\varphi\).  The
prefix through \(y_{n-1}\) is elementary by
\noderef{TashTreeProperPrefixElementary}.  By
\noderef{ColoredVertexSetElementary}, non-elementariness of the full set
gives two distinct displayed vertices with a common missing color
\(\beta\), where missingness has the meaning in
\noderef{MissingColors}.  Both vertices cannot lie in the elementary
prefix, and \noderef{TreeSequencePrefix} leaves only \(y_n\) outside it.
Thus \(\beta\) is missing at \(y_n\) and at \(y_j\) for some \(j<n\).

By \noderef{TashTreeFourUnusedColors}, the prefix through \(y_{n-2}\)
has four colors in \noderef{TreeFreeColors}.  The definitions
\noderef{TreeMissingColors}, \noderef{MissingColors}, and
\noderef{TreeUsedColors} say that each has a missing witness in that
prefix and is unused on its displayed colored edges.  The two remaining
exact identities \(e_{n-1},e_n\) use at most two colors.  Hence at least
two of the four colors are unused on all displayed colored edges of
\(T\).  Choose one, \(\alpha\ne\beta\), and a missing witness \(y_m\)
with \(m\le n-2\).

First suppose \(j=n-1\).  The color \(\beta\) is unused through
\(y_{n-1}\).  Indeed, if one of the relevant colored edges before
\(y_{n-1}\) had color \(\beta\), then \noderef{TashkinovTree} would
supply a missing-\(\beta\) witness at a still earlier vertex.  That
witness is distinct from \(y_{n-1}\) by vertex injectivity in
\noderef{TreeSequence}; by \noderef{MissingColors}, it and
\(y_{n-1}\) would be two distinct prefix vertices missing \(\beta\),
contrary to \noderef{ColoredVertexSetElementary}.  The edge \(e_n\)
has endpoint \(y_n\) by \noderef{TreeSequence}, so
\noderef{MissingColors} says it cannot have color \(\beta\).  Thus
\(\beta\) is unused on the whole tree.  Apply
\noderef{TashTreeControlledInterchange} to \(\alpha,\beta\) at
\(y_m,y_{n-1}\).  Its component is a path with those two controlled
ends and the returned coloring is again a minimum bad pair.  To apply
\noderef{MissingColorIncidenceEndpoint} at \(y_n\), its three premises
are explicit: \(\alpha\ne\beta\) by choice, \(\varphi\) is proper by
\noderef{KTriple}, and \(\beta\) is missing at \(y_n\).  Applying the
lemma with the two colors in the reversed order gives degree at most one
for the vertex-node \(y_n\).  If \(y_n\) lay in the controlled
component, the spanning path in
\noderef{TwoColorComponentSpanningPath} would therefore make it an end.
But its two ends are already \(y_m,y_{n-1}\), and all three vertices are
distinct by \noderef{TreeSequence}.  Hence \(y_n\) is outside the
component, and the outside-vertex clause of
\noderef{TwoColorComponentSwapMissingBehavior} shows that the swap
leaves the missing-\(\beta\) status at \(y_n\); the controlled-end
conclusion of \noderef{TashTreeControlledInterchange} makes \(\beta\)
missing at \(y_m\).  The controlled interchange also fixes
\(e_2,\ldots,e_{n-1}\).  The last edge-node is outside the component as
well.  If its color is \(\alpha\) or \(\beta\), its incidence with
\(y_n\) from \noderef{TreeSequence} is an adjacency in
\noderef{TwoColorIncidenceGraph}; membership of the edge-node would then
put \(y_n\) in the same connected support by
\noderef{TwoColorComponent}.  If its color is outside the pair or it is
uncolored, \noderef{TwoColorIncidenceGraph} gives it no adjacency at all,
so its isolated edge-node cannot lie in the component rooted at the
distinct vertex-node \(y_{n-1}\).  In either case the edge-node is
outside, and \noderef{TwoColorComponentSwap} fixes \(e_n\).  The same
definition fixes the uncolored \(e_1\).  Thus every displayed edge value
is fixed and \(\beta\), which was unused under \(\varphi\), is still
unused after the swap.
Taking \(j=m\) gives the conclusion.

Now suppose \(j\le n-2\).  If \(\beta\notin
\noderef{TreeUsedColors}(\varphi,T)\), take \(\varphi_1=\varphi\).
Otherwise let \(Q\) be the \((\alpha,\beta)\)-component rooted at
\(y_n\).  If \(Q\) meets the prefix through \(y_{n-2}\), then
\noderef{TashTreeLastIndexIntersectionShortening} applies: both colors
are missing on that prefix, \(y_n\) is outside it by vertex injectivity,
and its length is exactly \(n-2\).  This gives a contradiction.

Thus \(Q\) avoids every vertex of the prefix.  We claim that it also
avoids every displayed exact tree-edge node.  The edges through
\(e_{n-2}\) have both ends in the prefix, and \(e_{n-1}\) has its earlier
end there, by \noderef{TreeSequence}.  For the final edge, if its earlier
attachment were \(y_{n-1}\), then its ends would be
\(y_{n-1},y_n\).  Since \(n\ge3\), the defining condition in
\noderef{TashTreeSuffixPath} would then hold at \(n-1\), and the
leastness clause of \noderef{TashTreePathIndex} would give
\(\rho(T)\le n-1\), contrary to \(\rho(T)=n\).  Thus the earlier
attachment supplied by \noderef{TreeSequence} lies in the prefix through
\(y_{n-2}\).  Consequently every displayed edge has an endpoint in that
prefix.

Fix any displayed exact edge-node.  If its color is \(\alpha\) or
\(\beta\), then \noderef{TwoColorIncidenceGraph} joins it to the
prefix endpoint just identified.  If the edge-node belonged to \(Q\),
that adjacency and the connected-component meaning of
\noderef{TwoColorComponent} would put the endpoint in \(Q\), contrary to
vertex avoidance.  If the edge is uncolored or has a color outside
\(\{\alpha,\beta\}\), \noderef{TwoColorIncidenceGraph} gives its
edge-node no neighbors.  It is therefore an isolated component and,
since \(Q\) is the component rooted at the distinct vertex-node \(y_n\),
\noderef{TwoColorComponent} again excludes it from \(Q\).  This proves
avoidance of every displayed exact edge-node, including the isolated
off-pair case.

The color \(\alpha\) is unused on the whole displayed tree by its
choice above.  Apply \noderef{TashTreeDisjointComponentNormalization}
with retained-prefix cutoff \(p=n-2\), the fixed missing witnesses
\(y_m,y_j\), and the vertex- and edge-node avoidance just proved.  It
returns the explicit local component-swap coloring, proves it is again a
minimum bad pair, and makes \(\alpha\) missing at both \(y_m,y_n\).
Every displayed exact edge-node avoids \(Q\), so
\noderef{TwoColorComponentSwap} fixes every tree-edge value.  Since
\(\alpha\) was unused on the original tree, it remains outside
\noderef{TreeUsedColors} after the swap.  Rename this newly repeated
color \(\alpha\) as the output \(\beta\), and take \(j=m\).  This uses
only the local swap; no global transposition of color names is needed.
\end{proof}
